% ch8_c (书页 349-361 / p364-p376)
%--B-TAIL--
% ==== tail：8.5 节（红移与距离）收尾，承接书页 348 末的 (8.125)，
% ==== 至书页 349 中 8.6 节标题之前；置于本文件第一个 \section 之前 ====

% ---- 书页 349 / p364 ----
%JOIN%其中 $R$ 是天体的固有大小，$\theta$ 是观测到的角直径。在这两种情形下，我们都能导出与式 (8.123) 类似的公式；所幸，对宇宙学参数的那种难以驾驭的依赖关系为所有距离量度所共有，最后剩下的只是对红移的简单依赖：
\begin{equation*}
d_L = (1+z)d_M = (1+z)^{2}d_A,
\tag{8.126}
\end{equation*}
鼓励各位动手验证一下。这样，只要测得了其中一种距离，我们就能方便地换算成其他任何一种；或者，我们也可以独立地测量不同的距离，并利用式 (8.126) 来检验 RW 框架的自洽性。

趁我们还在思量距离，不妨也考虑一下：从红移为 $z$ 的天体发出的光，从发射时刻到现在流逝了多长时间。如果宇宙今天的年龄是 $t_0$，而光子被发射时宇宙的年龄是 $t_{*}$，那么\textbf{回溯时间}（lookback time）为
\begin{align*}
t_0 - t_{*} &= \int_{t_{*}}^{t_0}\mathrm{d}t\\
&= \int_{a_{*}}^{1}\frac{\mathrm{d}a}{aH(a)}\\
&= H_0^{-1}\int_0^{z_{*}}\frac{\mathrm{d}z'}{(1+z')E(z')}.
\tag{8.127}
\end{align*}
例如，考虑一个平直（$k=0$）、物质主导（$\rho=\rho_{\mathrm{M}}=\rho_{\mathrm{M}0}a^{-3}$）的宇宙。那么
\begin{equation*}
E(z) = (1+z)^{3/2},
\tag{8.128}
\end{equation*}
于是
\begin{align*}
t_0 - t_{*} &= H_0^{-1}\int_0^{z_{*}}\frac{\mathrm{d}z'}{(1+z')^{5/2}}\\
&= \frac{2}{3}H_0^{-1}\left[1-(1+z_{*})^{-3/2}\right].
\tag{8.129}
\end{align*}
令 $t_{*}\to 0$（$z_{*}\to\infty$），便得到物质主导宇宙的总年龄：
\begin{equation*}
t_0(\mathrm{MD}) = \frac{2}{3}H_0^{-1}.
\tag{8.130}
\end{equation*}
对于并非完全物质主导的宇宙，$\frac{2}{3}$ 这个因子不会完全正确，但对于合理的宇宙学参数取值，我们通常得到 $t_0\sim H_0^{-1}$。

\section{引力透镜}

在第 7 章中我们介绍了引力透镜（gravitational lensing）的概念：光在 Newton 引力场中的偏折与时间延迟。除了在太阳系内为广义相对论提供检验之外，透镜效应还出现在众多天体物理情境之中，并已成为现代宇宙学不可或缺的组成部分。\footnote{关于引力透镜的一篇优秀综述——我们的讨论即从中借鉴——参见 R.~Narayan and M.~Bartelmann, ``Lectures on Gravitational Lensing,'' 13th Jerusalem Winter School in Theoretical Physics, \texttt{http://arxiv.org/astro-ph/9606001}。}

% ---- 书页 350 / p365 ----
%FIG{8.5}: {引力透镜的几何，概括于透镜方程 (8.132) 之中。透镜的作用，是把平直 Minkowski 背景中本应观测到的角 $\beta$ 扭变为角 $\theta$。}

宇宙学透镜效应与我们先前讨论的情形有两个重要区别：Robertson--Walker 度规取代了 Minkowski 背景，而且透镜本身往往也比简单的点质量复杂。图 8.5 描绘了一种典型的透镜几何。在这一节的讨论中，我们将始终假设透镜是“薄”的——即其空间延展尺度远小于源、透镜与观者之间的距离。在这种情形下，我们可以合理地谈论一个唯一的到透镜的距离 $d_L$，以及透镜到源的距离 $d_{LS}$。

对于天空上一幅（可能很复杂的）像，我们用像的不同成分之间的一组夹角来描述它。这些夹角可以看作天空上的二维矢量。透镜的作用，是把没有任何偏折时本应观测到的那些角——例如源与透镜之间的夹角 $\vec{\beta}$——扭变为一幅由一组角 $\vec{\theta}$ 刻画的新像。我们全程假设这些角都很小。这一映射由\textbf{约化透镜角}（reduced lensing angle）$\vec{\alpha}=\vec{\theta}-\vec{\beta}$ 描述。按照图 8.5 所示的几何，它与真实偏转角 $\widehat{\vec{\alpha}}$ 之间的关系为
\begin{equation*}
\vec{\alpha} = \frac{d_{LS}}{d_S}\widehat{\vec{\alpha}}.
\tag{8.131}
\end{equation*}

% ---- 书页 351 / p366 ----
于是我们得到\textbf{透镜方程}（lens equation）
\begin{equation*}
\boxed{\,\vec{\beta} = \vec{\theta} - \frac{d_{LS}}{d_S}\widehat{\vec{\alpha}}.\,}
\tag{8.132}
\end{equation*}
透镜方程所描述的，无非是受扰时空中的光线追踪。

当然，我们应当仔细琢磨图中画出的那些“距离” $d_i$。透镜效应发生在膨胀的宇宙中，而宇宙还可能带有空间曲率。不过，只要我们\emph{定义}距离 $d_i$ 使得透镜方程所描述的几何关系得以成立，透镜方程就依然成立。换言之，这些距离是：在静态欧几里得空间背景中，给定夹角与横向物理尺度时我们本会推断出的距离。而这恰恰就是角直径距离 (8.125) 的定义。因此，本节中的所有距离都取为角直径距离。注意，角直径距离未必能够相加，故 $d_S\neq d_L+d_{LS}$。

举一个简单的例子：点质量透镜。在第 7 章研究 Newton 极限时我们曾经发现，光子穿过引力势 $\Phi$ 时的偏转角由下式给出：
\begin{equation*}
\widehat{\vec{\alpha}} = 2\int \vec{\nabla}_{\perp}\Phi\,\mathrm{d}s,
\tag{8.133}
\end{equation*}
对于碰撞参数为 $b$ 的点质量 $M$，上式化为
\begin{equation*}
\widehat{\vec{\alpha}} = \frac{4GM}{b}.
\tag{8.134}
\end{equation*}
碰撞参数可以写成 $b=d_L\theta$。透镜方程 (8.132) 便成为
\begin{equation*}
\beta = \theta - \frac{d_{LS}}{d_S d_L}\frac{4GM}{\theta}.
\tag{8.135}
\end{equation*}
富有启发性的做法是考虑最简单的情形：源与透镜共线（$\beta=0$）。此时源将被透镜成一个环绕透镜的\textbf{Einstein 环}（Einstein ring），其角距离由\textbf{Einstein 角}（Einstein angle）给出：
\begin{equation*}
\theta_{\mathrm{E}} = \sqrt{\frac{4GMd_{LS}}{d_L d_S}}.
\tag{8.136}
\end{equation*}
即便在更复杂的位形中，Einstein 角也为透镜效应设定了一个特征尺度。我们还可以定义一个与之相联系的距离尺度，即\textbf{Einstein 半径}（Einstein radius）：
\begin{equation*}
R_{\mathrm{E}} = \sqrt{\frac{4GMd_L d_{LS}}{d_S}}.
\tag{8.137}
\end{equation*}

% ---- 书页 352 / p367 ----
换算成厘米或其他物理单位时，别忘了在我们的所有方程中都有 $c=1$。为了对典型天体物理情形下透镜效应的大小有些感觉，我们可以考虑两种常见的场合：我们银河系内近似太阳质量的天体造成的“微透镜”（microlensing），以及星系或星系团造成的宇宙学透镜。前一情形下 Einstein 角将是毫角秒的量级，而后一情形下则是角秒的量级：
\begin{align*}
\theta_{\mathrm{E}} &= 0.9\sqrt{\left(\frac{M}{M_{\odot}}\right)\left(\frac{10\ \mathrm{kpc}}{D}\right)}\ \text{毫角秒}\\
&= 0.9\sqrt{\left(\frac{M}{10^{11}M_{\odot}}\right)\left(\frac{\mathrm{Gpc}}{D}\right)}\ \text{角秒}.
\tag{8.138}
\end{align*}
我们暂时继续讨论点质量透镜。尽管已经观测到若干例壮观的 Einstein 环，但多数时候我们不会幸运到源与透镜恰好完美对齐。这时我们可以求解 (8.135)，得到像角的两个数值：
\begin{equation*}
\theta_{\pm} = \frac{1}{2}\left(\beta \pm \sqrt{\beta^{2}+4\theta_{\mathrm{E}}^{2}}\right).
\tag{8.139}
\end{equation*}
位于 $\theta_{+}$ 处的像总在 Einstein 角之外，而 $\theta_{-}$ 处的像则在其内。实际上这个公式有些误导性：像的数目永远会是奇数；对点质量透镜而言，第三个像将位于与透镜自身相同的位置。

现在来考虑比点质量更一般的透镜。我们知道，偏转角可以借助 Newton 引力势由式 (8.133) 给出。我们可以沿从观者出发指向过去的测地线路径作积分，定义\textbf{透镜势}（lensing potential）：
\begin{equation*}
\psi(\vec{\theta}) = 2\frac{d_{LS}}{d_L d_S}\int \Phi(d_L\vec{\theta},s)\,\mathrm{d}s.
\tag{8.140}
\end{equation*}
利用透镜势，只需取梯度便可直接导出约化透镜角：
\begin{align*}
\vec{\alpha} &= \vec{\nabla}_{\theta}\psi\\
&= 2\frac{d_{LS}}{d_S}\int \vec{\nabla}_{\perp}\Phi\,\mathrm{d}s.
\tag{8.141}
\end{align*}
注意，角梯度 $\vec{\nabla}_{\theta}$ 与 $\vec{\nabla}_{\perp}$（对透镜所在处横向距离的梯度）相差一个因子 $d_L$。薄透镜近似使我们可以把积分收缩为在透镜位置上取值的量。我们还可以对透镜势作用（二维）拉普拉斯算符，得到\textbf{会聚}（convergence）$\kappa$：
\begin{align*}
\kappa(\vec{\theta}) &\equiv \frac{1}{2}\nabla_{\theta}^{2}\psi\\
&= \frac{d_L d_{LS}}{d_S}\int \nabla^{2}\Phi\,\mathrm{d}s.
\tag{8.142}
\end{align*}

% ---- 书页 353 / p368 ----
会聚可以被看作对积分质量密度的一种量度。我们可以把上面的表达式反过来，用会聚同时写出透镜势与约化偏转角：
\begin{equation*}
\psi(\vec{\theta}) = \frac{1}{\pi}\int \kappa(\vec{\theta}\,')\ln|\vec{\theta}-\vec{\theta}\,'|\,\mathrm{d}^{2}\theta'
\tag{8.143}
\end{equation*}
以及
\begin{equation*}
\vec{\alpha}(\vec{\theta}) = \frac{1}{\pi}\int \kappa(\vec{\theta}\,')\frac{\vec{\theta}-\vec{\theta}\,'}{|\vec{\theta}-\vec{\theta}\,'|}\,\mathrm{d}^{2}\theta'.
\tag{8.144}
\end{equation*}
要检验这些方程，请记住这些矢量只定义在两个横向维度上。

会聚描述的是引力透镜对光线的聚焦。这种聚焦使源显得更大（就像放大镜一样）。根据 Liouville 定理——源所发射光子的相空间密度守恒——源的面亮度在透镜作用下保持不变；因此尺度的增大导致亮度的放大。与此同时，光线穿过透镜时的扭转还会引起畸变，导致像的形状发生切变。为了描述这两种现象，我们考虑透镜映射的导数所构成的 $2\times 2$ 矩阵
\begin{equation*}
A_{ij} \equiv \frac{\partial\beta^{i}}{\partial\theta^{j}}.
\tag{8.145}
\end{equation*}
注意，上、下指标之间实际上没有区别，因为它们定义在二维欧几里得平面上。由于 $\vec{\beta}=\vec{\theta}-\vec{\alpha}$，我们有
\begin{align*}
A_{ij} &= \delta_{ij} - \frac{\partial\alpha^{i}}{\partial\theta^{j}}\\
&= \delta_{ij} - \psi_{ij},
\tag{8.146}
\end{align*}
这里我们引入了记号
\begin{equation*}
\psi_{ij} \equiv \frac{\partial^{2}\psi}{\partial\theta^{i}\partial\theta^{j}}.
\tag{8.147}
\end{equation*}
矩阵 $A$ 编码了透镜映射的局域性质。它的逆矩阵被称为\emph{放大张量}（magnification tensor）
\begin{equation*}
\boldsymbol{M} = \frac{\partial\vec{\theta}}{\partial\vec{\beta}} = A^{-1}.
\tag{8.148}
\end{equation*}

% ---- 书页 354 / p369 ----
这个名字从何而来？透镜把由 $\vec{\beta}$ 描述的面元扭变为由 $\vec{\theta}$ 描述的面元，面积的改变由这个映射的雅可比行列式描述，而后者无非是 $\boldsymbol{M}$ 的行列式。这个行列式被定义为\textbf{放大率}（magnification）$\mu$：
\begin{equation*}
\mu = |\boldsymbol{M}| = \frac{1}{|A|}.
\tag{8.149}
\end{equation*}
$\mu$ 的绝对值告诉我们源的亮度实际改变了多少；$\mu$ 可能为负，这意味着像的奇偶性被翻转了。我们说“放大”，是因为只有当透镜和源在天空中彼此靠近时透镜效应才会惹人注意，此时聚焦效应只会带来视亮度的增加；而一个远离源（在天空中位置相距甚远）的透镜所导致的光度极微小的减少，则永远不会被注意到。（如果存在多重像，所有像的亮度之和将超过未畸变源的亮度。）

$A$ 的分量可以分解为会聚与切变的效应。对于会聚，由 $\kappa=\frac{1}{2}\nabla_{\theta}^{2}\psi$ 可得
\begin{equation*}
\kappa = \frac{1}{2}(\psi_{11}+\psi_{22}).
\tag{8.150}
\end{equation*}
而\textbf{切变}（shear）则使源的形状发生畸变；如果初始为圆形的源被畸变为椭圆度为 $\gamma$、方位角为 $\phi$ 的椭圆，我们定义切变的两个分量为
\begin{align*}
\gamma_1 &= \gamma\cos(2\phi)\\
\gamma_2 &= \gamma\sin(2\phi),
\tag{8.151}
\end{align*}
于是总切变为 $\gamma=\sqrt{\gamma_1^{2}+\gamma_2^{2}}$。用透镜势表出，这些分量为
\begin{align*}
\gamma_1 &= \frac{1}{2}(\psi_{11}-\psi_{22})\\
\gamma_2 &= \psi_{12}=\psi_{21}.
\tag{8.152}
\end{align*}
反解这些关系以求出 $A$ 的分量，得到
\begin{equation*}
A = \begin{pmatrix} 1-\kappa-\gamma_1 & -\gamma_2\\ -\gamma_2 & 1-\kappa+\gamma_1 \end{pmatrix}.
\tag{8.153}
\end{equation*}
于是我们可以用会聚与切变来表示放大率：
\begin{equation*}
\mu = \frac{1}{(1-\kappa)^{2}-\gamma^{2}}.
\tag{8.154}
\end{equation*}

% ---- 书页 355 / p370 ----
透镜效应的这些特征在观测宇宙学中正变得日益重要。显然令人感兴趣的是所谓“强透镜”（strong lensing）的情形：源位于透镜的 Einstein 半径之内，这时可能形成多重像。通过观测同一源的多个像，我们可以推断透镜质量分布的性质（例如用来搜寻暗物质）；我们还可以利用不同路径上的时间延迟来测量 Hubble 常数，并利用透镜事件的统计频率来约束其他宇宙学参数。不过，透镜效应不必很强也能产生重要影响。“弱透镜”（weak lensing）——源与透镜相距超过一个 Einstein 半径——一般导致少量的放大与切变，若不预先知道源的性质便无法探测。然而，切变效应可以用统计方法探测：考察成千上万个星系的形状（假定它们的取向本质上是随机的）。弱透镜切变导致这些形状出现相关联的畸变，从而可以揭示观者与遥远源之间物质分布的大量信息。

\section{我们的宇宙}

在讨论 FRW 宇宙学的行为时，我们屡次提及与我们身处的宇宙相对应的宇宙学参数的实际取值。现在让我们更系统一些：既讨论我们今天看到的宇宙，也讨论向早期时刻的合理外推。我们的讨论必然简略，这既出于篇幅考虑，也因为宇宙学是一个活跃的研究领域；想要了解当前观点的最新描述，请查阅近期的综述文章。

我们对膨胀率的许多直接测定，依靠的是把光度距离公式 (8.123) 应用于某种其内禀光度被假定已知的天体，这类天体称为\textbf{标准烛光}（standard candles）。（偶尔我们也测量其内禀尺度被假定已知的物体的角直径：标准尺［standard rulers］。）例如，Hubble 常数就是用各种标准烛光测量的，不同方法达成的共识收敛于前面提到的值 $H_0=70\pm 10$ km/sec/Mpc。高红移处对线性 Hubble 定律 (8.109) 的偏离可以提供关于密度参数 $\Omega_{i0}$ 的信息，但前提是我们必须拥有非常明亮、且内禀光度已被精确掌握的天体。Ia 型超新星（Type Ia supernovae）恰好提供了这样的天体；它们被认为是白矮星吸积了足够多的质量、以至于超过 Chandrasekhar 极限之后发生的爆发。由于 Chandrasekhar 极限接近普适，相应的爆发在本质上亮度相同（而且一部分内禀变化实际上可以通过追踪亮度随时间的演化而加以扣除）。正是在红移 $z>0.3$ 处对 Ia 型超新星的测量，提供了宇宙学常数非零的第一个直接证据；这些观测意味着 $\Omega_{\Lambda}$ 实际上大于 $\Omega_{\mathrm{M}}$。回想一下：物质是无压强的，$p_{\mathrm{M}}=0$，而真空能则伴随着负压强，$p_{\Lambda}=-\rho_{\Lambda}$。代入第二个弗里德曼方程 (8.68)，我们发现一个同时含有物质与 $\Lambda$ 的宇宙服从

% ---- 书页 356 / p371 ----
\begin{equation*}
\frac{\ddot{a}}{a} = -\frac{4\pi G}{3}(\rho_{\mathrm{M}} - 2\rho_{\Lambda}).
\tag{8.155}
\end{equation*}
于是，如果 $\rho_{\Lambda}$ 相对于 $\rho_{\mathrm{M}}$ 足够大（正如超新星观测所表明的），我们就可以有 $\ddot{a}>0$：一个加速膨胀的宇宙（其含义如 8.4 节所述）。

物质密度本身的测量方法多种多样，通常包括通过寻找成团物质的引力效应来测量密度 $\rho_{\mathrm{M}}$，然后再外推到大尺度。由于 $\rho_{\mathrm{M}}=(3H^{2}/8\pi G)\Omega_{\mathrm{M}}$，用这种方式得到的限值常常以 $\Omega_{\mathrm{M}}h^{2}$ 的形式给出，其中 $h$ 定义于式 (8.70)。如今 $H_0$ 的不确定度看来已足够小，取 $h^{2}\approx 0.5$ 相当稳妥，从此以后我们就这么做。当代的大多数方法都与如下结果一致：
\begin{equation*}
\Omega_{\mathrm{M}0} = 0.3\pm 0.1.
\tag{8.156}
\end{equation*}
在宇宙学常数获得充分的观测证据之前，如此低的物质密度有时曾被当作空间负弯曲（$\kappa<0$）的迹象。

除物质与宇宙学常数之外，宇宙中还有辐射。普通光子是辐射密度中最显然的成分，但任何相对论性粒子都会作出贡献。就光子而言，绝大部分能量密度蕴藏于宇宙微波背景——大爆炸遗留下来的辐射——之中。除光子外，辐射成分唯一显然的候选者是中微子。我们预期遗迹背景中微子的数密度与光子相当；光子的数密度可能还要稍大一些，因为在中微子数目固定下来之后，光子仍能继续产生。然而，如果中微子的质量足够大（大于约 $10^{-4}$ eV），它们今天将已经成为非相对论性的，从而对物质而非辐射作出贡献。目前关于中微子质量的看法提示情况很可能如此，但也并非完全清楚。此外，可以设想在我们已知的粒子之外，还存在迄今尚未探测到的无质量粒子（尽管它们不能太多，否则会抑制大尺度结构的形成）。总而言之，总辐射密度看来很可能与光子密度处于同一数量级；在这种情形下我们将有
\begin{equation*}
\Omega_{\mathrm{R}0} \sim 10^{-4}.
\tag{8.157}
\end{equation*}
如前所述，辐射密度低于物质密度并不令人惊讶，因为随着宇宙的膨胀，前者衰减得更快。辐射密度按 $a^{-4}$ 变化，而物质中的密度按 $a^{-3}$ 变化；因此物质-辐射相等发生的红移为
\begin{equation*}
z_{\mathrm{eq}} \approx \frac{\Omega_{\mathrm{M}0}}{\Omega_{\mathrm{R}0}} \sim 3\times 10^{3}.
\tag{8.158}
\end{equation*}

% ---- 书页 357 / p372 ----
对宇宙学参数的另一个关键约束来自微波背景温度的各向异性。平均温度为 $T_{\mathrm{CMB}}=2.74\,\mathrm{K}$，但 1992 年 COBE 卫星发现了随地点起伏、幅度为 $\Delta T/T\sim 10^{-5}$ 的涨落。这些各向异性有多种来源，其中包括复合时光子移出势阱带来的引力红移/蓝移（Sachs--Wolfe 效应，在大角尺度上占主导）、最后散射面上的内禀温度涨落（在小角尺度上占主导），以及等离子体运动的多普勒效应。描写 CMB 各向异性演化的物理超出了本书的范围。整幅天空的 CMB 温度图显然包含大量信息，但没有任何理论预言任意给定点处的温度应该是多少；现代理论预言的，一般是任意给定角尺度上各向异性大小的期望值。因此，我们把各向异性场按球谐函数展开：
\begin{equation*}
\frac{\Delta T}{T}(\theta,\phi) = \sum_{lm} a_{lm}Y_{lm}(\theta,\phi).
\tag{8.159}
\end{equation*}
$|a_{lm}|^{2}$ 的期望值很可能与 $m$ 无关；否则各向异性的统计特性会随天空中的位置而改变（尽管我们对此应当保持开放的心态）。因此，有待测量的相关参数是
\begin{equation*}
C_l = \langle|a_{lm}|^{2}\rangle.
\tag{8.160}
\end{equation*}
由于对任意固定的 $l$，$m$ 有 $2l+1$ 个可能的取值（从 $-l$ 到 $l$），除了最低的几个 $l$ 之外，对 $a_{lm}$ 的独立测量足以精确确定它们的期望值。在非常小的 $l$ 处这种不可约的不确定度被称为宇宙方差（cosmic variance）。

已有大量实验测量了 $C_l$（所谓的 CMB 功率谱），而改进这些测量在今后若干年内很可能仍是一项重要任务。（除温度各向异性之外，CMB 的偏振也包含大量信息，是另一个需要相当多实验努力的目标。）要把这些观测转化为有用的信息，我们需要一个具体的理论来预言 CMB 功率谱如何作为宇宙学参数的函数而变化。有两种领先的可能性（尽管其中一种比另一种要领先得多）：要么密度扰动在极早的时刻就被印记在所有尺度上——甚至包括物理波长 $\lambda$ 远大于 Hubble 半径 $H^{-1}$ 的模式；要么局部的动力学机制在所有时期充当各向异性的源。后一种可能性已基本上被 CMB 数据排除：如果各向异性是连续产生的，我们预期会看到相当平滑、无特征的 $C_l$ 谱，而观测却显示出显著的结构。因此，更流行的做法是设想某种原初的扰动源，例如暴胀（下一节讨论）。暴胀扰动是绝热的——物质密度中的扰动与辐射密度中的扰动相关联——并且在所有尺度上幅度几乎相等。有了这个输入，我们就可以对 $C_l$ 作出作为所有宇宙学参数函数的明确预言。迄今从实验得到的最重要约束或许是：宇宙在空间上是平直的，或至少非常接近平直；$|\Omega_{\mathrm{c}0}|<0.1$。结合物质密度的测量 $\Omega_{\mathrm{M}}\approx 0.3$，我们得出结论：真空能密度参数应为

% ---- 书页 358 / p373 ----
\begin{equation*}
\Omega_{\Lambda 0} = 0.7\pm 0.1.
\tag{8.161}
\end{equation*}
这与上面描述的 Ia 型超新星结果符合得很好；这里描述的调和图像（concordance picture）正是图 8.4 所指示的那一种。利用 $H_0=70$ km/sec/Mpc 把密度参数换算成物理能量密度，得到
\begin{equation*}
\rho_{\mathrm{vac}} \approx 10^{-8}\ \mathrm{erg/cm^{3}},
\tag{8.162}
\end{equation*}
如 4.5 节讨论真空能时所述。

我们的现今宇宙概貌还有一个值得注意的特征。前面提到，我们宇宙中约 $30\%$ 的能量密度由物质组成。但对宇宙学家而言，“物质”是任何非相对论性粒子的集合；我们通过其引力影响推断出的物质，未必是我们在地球上的经验所熟悉的那种普通物质。\textbf{普通物质}（ordinary matter）指由原子及其成分（质子、中子和电子）构成的一切东西；这包括宇宙中所有的恒星、行星、气体和尘埃，无论能否直接看到。这类物质偶尔也被称为\emph{重子物质}（baryonic matter），其中重子包括质子、中子以及相关的粒子（携带一个被称为重子数的守恒量子数的强相互作用粒子）。当然，电子在概念上是普通物质的重要部分，但按质量计算，与质子和中子相比可以忽略：
\begin{gather*}
m_p = 0.938\ \mathrm{GeV}\\
m_n = 0.940\ \mathrm{GeV}\\
m_e = 0.511\times 10^{-3}\ \mathrm{GeV}.
\tag{8.163}
\end{gather*}
换句话说，普通物质的质量绝大部分来自重子。

事实证明，普通的重子物质远不足以解释观测到的密度 $\Omega_{\mathrm{M}}\approx 0.3$。我们目前对重子密度的最佳估计给出
\begin{equation*}
\Omega_{\mathrm{b}} = 0.04\pm 0.02,
\tag{8.164}
\end{equation*}
这些误差限按大多数标准衡量都是保守的。这一测定来自多种方法：直接计数重子（最不精确的方法）、与 CMB 功率谱（上文讨论过）的一致性，以及与大爆炸核合成（下文讨论）对轻元素丰度的预言相符。因此，物质密度的大部分必定采取\textbf{非重子暗物质}（nonbaryonic dark matter）的形式，我们将把它简称为“暗物质”。（重子也可以是暗的，但越来越普遍的用法是把这一术语保留给非重子成分。）粒子物理标准模型中的每一种已知粒子，作为这种暗物质的候选者都已被排除。幸运的是，标准模型之外还有若干看似合理的候选者，包括中性伴随子（neutralino，超对称性所预言的附加稳定粒子中最轻者，质量 $\geq 100$ GeV）和轴子（axion，因假设的 Peccei--Quinn 对称性自发破缺而产生的轻赝标量粒子，引入该对称性是为了解释强相互作用中 CP 的守恒，质量 $\sim 10^{-4}$ eV）。关于暗物质，我们仅有的几点确切认识之一是：它必须是冷的——不仅今天是非相对论性的，而且在很长一段时间里都必须如此。如果暗物质是热的，它就会自由流出高密度区，抑制星系的形成。关于冷暗物质（cold dark matter，CDM）我们知道的另一件事是，它与普通物质的相互作用必须非常微弱，这样才能解释它为何迄今未被探测到。尽管如此，环境中的暗物质粒子偶尔仍可能与地球实验室里经过仔细屏蔽的探测器发生散射；通过搜寻这种散射的效应来直接探测暗物质，将是今后若干年里另一项重要的实验努力。

% ---- 书页 359 / p374 ----
$\Omega_{\mathrm{M}}=0.3$、$\Omega_{\Lambda}=0.7$ 的这幅图像，似乎能拟合种类多得令人印象深刻的观测数据。这幅图中最令人惊讶的部分是宇宙学常数。在第 4 章中我们提到，对真空能的朴素估计给出的结果比测量值大许多个数量级。实际上有三个相互关联的谜题：为什么宇宙学常数比我们预期的小这么多？构成当前宇宙 $70\%$ 的那份小小的非零能量，其起源是什么？还有，为什么真空能的当前值与物质密度处于同一数量级？最后一个问题尤其严峻，因为真空能与物质密度彼此之间演化得很快：
\begin{equation*}
\frac{\Omega_{\Lambda}}{\Omega_{\mathrm{M}}} \propto a^{3}.
\tag{8.165}
\end{equation*}
如果 $\Omega_{\mathrm{M}}$ 和 $\Omega_{\Lambda}$ 今天彼此相当，那么过去真空能会小得探测不到，而将来物质密度将变得可以忽略。这个“巧合问题”（coincidence problem）迄今仍完全是个谜。有人提出的解法涉及“人择原理”（anthropic principle）：如果宇宙存在许多不同的部分（在空间上，甚至在波函数的不同分支上），其中宇宙学常数取非常不同的数值，那么智慧生命最有可能出现在那些绝对值不太大的地方——一个大的正 $\Lambda$ 会在星系形成之前把粒子撕开，而一个大的负 $\Lambda$ 会使宇宙在生命演化出来之前就再度坍缩。对观测到的真空能的人择解释与数据拟合得很好，尽管为解释这一个量就需要求助如此精巧的方案，在一些人看来有点夸张。

另一种可能与巧合问题有关（也可能无关）的想法是：我们探测到的并非非零的宇宙学常数，而是一个紧密切拟真空能性质的动力学成分。对这种可能性的考虑，促使宇宙学家创造了\textbf{暗能量}（dark energy）这个术语，用来称呼无论实际上是什么的被探测到的东西——不管它是动力学的，还是最终被证明就是宇宙学常数。关于暗能量我们知道的是：它在空间中的分布相对平滑（否则它就会像暗物质那样，通过其局部引力场被探测到），并且随时间演化得很慢（否则它就不会像超新星数据所表明的那样使宇宙加速）。暗能量的一个简单的动力学候选者由缓慢滚动的标量场提供。考虑一个具有通常作用量的场 $\phi$：

% ---- 书页 360 / p375 ----
\begin{equation*}
S = \int \mathrm{d}^{4}x\sqrt{-g}\left[-\frac{1}{2}g^{\mu\nu}\nabla_{\mu}\phi\nabla_{\nu}\phi - V(\phi)\right],
\tag{8.166}
\end{equation*}
其能量-动量张量为
\begin{equation*}
T_{\mu\nu} = \nabla_{\mu}\phi\nabla_{\nu}\phi + \left[\frac{1}{2}g^{\rho\sigma}\nabla_{\rho}\phi\nabla_{\sigma}\phi - V(\phi)\right]g_{\mu\nu}
\tag{8.167}
\end{equation*}
而运动方程为
\begin{equation*}
\Box\phi - \frac{\mathrm{d}V}{\mathrm{d}\phi} = 0.
\tag{8.168}
\end{equation*}
假设场在空间中完全均匀（$\partial_i\phi=0$）。于是利用 Christoffel 符号 (8.44)，我们可以用时间导数和 Hubble 常数表示 d'Alembert 算符，把 (8.168) 写成
\begin{equation*}
\ddot{\phi} + 3H\dot{\phi} + \frac{\mathrm{d}V}{\mathrm{d}\phi} = 0.
\tag{8.169}
\end{equation*}
我们看到，Hubble 参量起到摩擦项的作用；场会趋于沿势滚落，但当 $H$ 太大时运动将被阻尼。因此，具有足够平坦的势的标量场（如图 8.6 所绘）将滚动得非常缓慢，导致动能远小于势能 $V(\phi)$。此时能量-动量张量为
\begin{equation*}
T_{\mu\nu} \approx -V(\phi)g_{\mu\nu},
\tag{8.170}
\end{equation*}
%FIG{8.6}: {缓慢滚动的标量场的势能。}
其中 $\phi\approx$ 常数。与式 (4.96) 比较，我们看到标量场势正在摹拟一种真空能。举一个简单的例子，考虑二次势 $V(\phi)=\frac{1}{2}m^{2}\phi^{2}$。这时 (8.169) 描写一个阻尼谐振子，若 $H>m$ 就会出现过阻尼。但在粒子物理单位制中，今天的 Hubble 常数为 $H_0\approx 10^{-33}$ eV，所以这个标量场的质量将不得不比式 (8.163) 中我们熟悉的基本粒子的质量小得难以置信。这看起来像是不自然的微调。尽管如此，动力学暗能量模型正在被积极探索，部分是出于这样的希望：它们能以某种方式通向巧合问题的解决。

% ---- 书页 361 / p376 ----
有了对当代状况的这一认识，我们就可以设想，早期宇宙必须是什么样子，才能产生出我们今天看到的一切。为了物理直观，标记所考虑的时代时，用温度来表示往往比用红移或大爆炸以来的时间更有帮助。今天的温度是
\begin{equation*}
T_0 = 2.74\ \mathrm{K} = 2.4\times 10^{-4}\ \mathrm{eV}.
\tag{8.171}
\end{equation*}
当然，这里说“温度”，我们指的是宇宙微波背景的表观黑体温度；实际上 CMB 自复合以来就不再处于热平衡，所以对这个概念不宜过分当真。在绝热膨胀下，随着每种相对论性粒子红移，温度都在下降，我们有 $T\propto a^{-1}$。但在早期宇宙的特定时刻会发生非绝热相变；在这种情形下温度实际上并不会升高，只是下降得更加缓慢。为了帮助建立温度、密度与尺度因子之间的关系，我们引入\textbf{相对论性自由度的有效数目}（effective number of relativistic degrees of freedom）的两种不同量度：$g_{*}$ 与 $g_{*S}$（其中 $S$ 代表熵）。考虑一组玻色子与费米子物种，每一种都有自己的有效温度 $T_i$ 和自旋态数目 $g_i$。例如，无质量光子有两个自旋态，故 $g_{\gamma}=2$；有质量的 $\frac{1}{2}$ 自旋费米子也有两个自旋态，故 $g_{e^-}=g_{e^+}=2$。相对论性自由度有效数目的这两个不同版本满足
\begin{equation*}
g_{*} = \sum_{\mathrm{bosons}} g_i\left(\frac{T_i}{T}\right)^{4} + \frac{7}{8}\sum_{\mathrm{fermions}} g_i\left(\frac{T_i}{T}\right)^{4}
\tag{8.172}
\end{equation*}
而
\begin{equation*}
g_{*S} = \sum_{\mathrm{bosons}} g_i\left(\frac{T_i}{T}\right)^{3} + \frac{7}{8}\sum_{\mathrm{fermions}} g_i\left(\frac{T_i}{T}\right)^{3}.
\tag{8.173}
\end{equation*}
神秘的 $\frac{7}{8}$ 因子，来源于计算平衡分布函数时玻色统计与费米统计的差别。对任何处于热平衡的物种，其温度 $T_i$ 将等于背景温度 $T$；但也可能有物种在更低的温度下就已退耦，它们对相对论性自由度有效数目的贡献较小。之所以需要定义两种不同的量度，是因为它们扮演的角色不同：第一个通过
\begin{equation*}
\rho_{\mathrm{R}} = \frac{\pi^{2}}{30}g_{*}T^{4},
\tag{8.174}
\end{equation*}
把温度与（相对论性物种的）能量密度联系起来，
而第二个则把温度与尺度因子联系起来：
\begin{equation*}
T \propto g_{*S}^{-1/3}a^{-1}.
\tag{8.175}
\end{equation*}
